% ============================================================
%  PROYECTO INTEGRADOR 1 — MAESTRÍA EN CIENCIA DE DATOS Y ANALÍTICA
%  Identificación y Predicción de Informalidad Laboral en Colombia
%  Juan Andrés Montoya · Julián David Mejía
% ============================================================
\documentclass[12pt, a4paper]{article}

% ----- Idioma y codificación -----
\usepackage[spanish,es-tabla]{babel}
\usepackage[utf8]{inputenc}
\usepackage[T1]{fontenc}

% ----- Diseño de página -----
\usepackage[top=3cm, bottom=2.5cm, left=3cm, right=2.5cm]{geometry}
\usepackage{setspace}
\onehalfspacing
\usepackage{parskip}
\setlength{\parindent}{0pt}
\setlength{\parskip}{6pt}

% ----- Tipografía -----
\usepackage{lmodern}
\usepackage{microtype}

% ----- Gráficos y colores -----
\usepackage{graphicx}
\usepackage{xcolor}
\definecolor{azulDane}{HTML}{1e3a5f}
\definecolor{azulClaro}{HTML}{2e6da4}
\definecolor{rojoRiesgo}{HTML}{C0392B}
\definecolor{grisCode}{HTML}{f5f5f5}
\definecolor{verdeCodigo}{HTML}{276221}
\definecolor{moradoCodigo}{HTML}{5c0099}

% ----- Tablas -----
\usepackage{booktabs}
\usepackage{tabularx}
\usepackage{longtable}
\usepackage{array}
\usepackage{multirow}
\usepackage{colortbl}
\newcolumntype{L}[1]{>{\raggedright\arraybackslash}p{#1}}
\newcolumntype{C}[1]{>{\centering\arraybackslash}p{#1}}
\newcolumntype{R}[1]{>{\raggedleft\arraybackslash}p{#1}}

% ----- Código fuente -----
\usepackage{listings}
\lstdefinestyle{pythonstyle}{
    backgroundcolor=\color{grisCode},
    commentstyle=\color{verdeCodigo}\itshape,
    keywordstyle=\color{moradoCodigo}\bfseries,
    stringstyle=\color{azulClaro},
    basicstyle=\ttfamily\small,
    breaklines=true,
    captionpos=b,
    keepspaces=true,
    numbers=left,
    numberstyle=\tiny\color{gray},
    numbersep=8pt,
    showspaces=false,
    showstringspaces=false,
    tabsize=4,
    frame=single,
    rulecolor=\color{gray!30},
    xleftmargin=12pt,
    framexleftmargin=12pt,
}
\lstset{style=pythonstyle, language=Python}

% ----- Encabezados y pies -----
\usepackage{fancyhdr}
\pagestyle{fancy}
\fancyhf{}
\fancyhead[L]{\small\textcolor{gray}{Montoya \& Mejía — Informalidad Laboral Colombia}}
\fancyhead[R]{\small\textcolor{gray}{PI-1 · Maestría CD\&A}}
\fancyfoot[C]{\thepage}
\renewcommand{\headrulewidth}{0.4pt}

% ----- Hipervínculos -----
\usepackage[colorlinks=true, linkcolor=azulDane, citecolor=azulDane,
            urlcolor=azulClaro, bookmarks=true]{hyperref}
\usepackage{url}

% ----- Matemáticas -----
\usepackage{amsmath, amssymb}

% ----- Figuras y flotantes -----
\usepackage{float}
\usepackage{caption}
\usepackage{subcaption}
\captionsetup{font=small, labelfont=bf}

% ----- Cajas destacadas -----
\usepackage{tcolorbox}
\tcbuselibrary{skins, breakable}
\newtcolorbox{hallazgo}[1][]{
  enhanced, breakable,
  colback=azulDane!5, colframe=azulDane,
  fonttitle=\bfseries,
  title={Hallazgo clave}, #1
}
\newtcolorbox{decision}[1][]{
  enhanced, breakable,
  colback=azulClaro!8, colframe=azulClaro,
  fonttitle=\bfseries,
  title={Decisión de diseño}, #1
}

% ----- Listas -----
\usepackage{enumitem}

% ----- Numeración de secciones -----
\setcounter{secnumdepth}{3}
\setcounter{tocdepth}{3}

% ============================================================
\begin{document}

% ============================================================
% PORTADA
% ============================================================
\begin{titlepage}
  \thispagestyle{empty}
  \begin{center}
    \vspace*{1cm}

    {\Large\bfseries\color{azulDane}
      Maestría en Ciencia de Datos y Analítica\\[0.3em]
      \normalsize Línea de énfasis: Ciencia de Datos y Analítica
    }

    \vspace{0.8cm}
    \rule{\textwidth}{1.5pt}
    \vspace{0.5cm}

    {\Huge\bfseries\color{azulDane}
      Identificación y Predicción de\\[0.3em]
      Informalidad Laboral en Colombia
    }

    \vspace{0.4cm}
    {\Large\color{azulClaro}
      Modelado predictivo sobre trabajadores independientes\\
      usando la GEIH 2024 (DANE)
    }

    \vspace{0.4cm}
    \rule{\textwidth}{0.4pt}

    \vspace{1.2cm}
    {\large\bfseries Proyecto Integrador 1}\\[0.3em]
    {\normalsize Materias: SI7009 Aprendizaje Automático · SI7006 Almacenamiento y
    Procesamiento de Grandes Datos · SI7007 Visualización de Datos}

    \vspace{1.5cm}
    \begin{tabular}{cc}
      \large\bfseries Juan Andrés Montoya & \large\bfseries Julián David Mejía\\[0.2em]
      \normalsize Maestrante CD\&A & \normalsize Maestrante CD\&A
    \end{tabular}

    \vfill

    \begin{tabular}{ll}
      \textbf{Fecha de entrega:} & Mayo 2026\\[0.2em]
      \textbf{Fuente de datos:} & DANE — GEIH 2024 (enero–diciembre)\\[0.2em]
      \textbf{Herramientas:} & Python 3.11 · LightGBM · Streamlit · DuckDB · Plotly\\[0.2em]
      \textbf{Despliegue:} & Streamlit Community Cloud (enlace público)\\
    \end{tabular}

    \vspace{0.8cm}
    {\small\textcolor{gray}{Fuente: DANE — Gran Encuesta Integrada de Hogares 2024.
    El DANE no avala los resultados del análisis.}}
  \end{center}
\end{titlepage}

% ============================================================
% RESUMEN
% ============================================================
\newpage
\thispagestyle{empty}
\begin{abstract}
\noindent Este trabajo construye un sistema de predicción de informalidad laboral para
trabajadores independientes (cuenta propia) en Colombia a partir de los microdatos de la
Gran Encuesta Integrada de Hogares 2024 (GEIH, DANE). La pregunta orientadora es:
\textit{¿Por qué los trabajadores independientes permanecen informales, y qué condiciones
deben cambiar para que la formalización sea una opción viable?}

Se entrena un modelo LightGBM sobre 351.721 registros (filtrados a 151.136 trabajadores
por cuenta propia), con 13 características socioeconómicas y laborales, y se alcanza un
\textbf{F1 = 0{,}9576} y \textbf{AUC-ROC = 0{,}9853} en el conjunto de prueba.
La interpretación SHAP revela que \textit{microempresa} (tamaño del establecimiento $\leq$10
personas) es el predictor más determinante (contribución SHAP media = 1{,}397), seguido
de los años de educación y la rama de actividad económica.
Los resultados se presentan en un dashboard Streamlit con cinco pestañas interactivas,
desplegado públicamente en Streamlit Community Cloud, que incluye un módulo de predicción
individual con recomendaciones de política pública generadas por inteligencia artificial.
\end{abstract}

\tableofcontents
\newpage

% ============================================================
\section{Introducción}
% ============================================================

La informalidad laboral es uno de los fenómenos estructurales más persistentes del
mercado de trabajo colombiano. Según el DANE, entre el 56\% y el 58\% de los ocupados
trabajan en condiciones de informalidad, entendida como ausencia de cotización activa al
sistema de seguridad social. Esta realidad afecta de forma desproporcionada a mujeres,
jóvenes, habitantes de zonas rurales y trabajadores con bajo nivel educativo.

A pesar de su magnitud, los instrumentos disponibles para el Ministerio del Trabajo y el
SENA se apoyan en estadísticas agregadas a nivel departamental o sectorial. El Banco
Interamericano de Desarrollo ha señalado que \textit{``es común implementar programas sin
diagnósticos comprehensivos de las brechas''} y que existen \textit{``grandes debilidades en
materia de focalización''} \cite{melendez2014}. Las evaluaciones de la Ley 1429 de 2010
(Ley de Formalización y Generación de Empleo) confirman que sus efectos fueron
heterogéneos entre subgrupos: en el corto plazo aumentó la probabilidad de empleo pero
también incrementó la informalidad en ciertos perfiles \cite{arizacedano2017}.

\textbf{Objetivo central.} Construir un modelo de clasificación binaria que estime la
probabilidad individual de informalidad laboral para trabajadores independientes
(cuenta propia) a partir de sus características socioeconómicas y laborales observables.
El resultado es un score de riesgo individual que permite a entidades como el Ministerio
del Trabajo, el SENA o las alcaldías priorizar intervenciones de formalización hacia los
perfiles con mayor probabilidad de permanecer informales.

\textbf{Estructura del documento.} La Sección~\ref{sec:marco} presenta el marco teórico y
los antecedentes. La Sección~\ref{sec:ml} describe el desarrollo metodológico completo
de la cadena de Machine Learning (EDA, ingeniería de características, modelado, evaluación).
La Sección~\ref{sec:tecnologia} explica las decisiones de ingeniería de datos y arquitectura.
La Sección~\ref{sec:viz} documenta las decisiones de visualización y comunicación. La
Sección~\ref{sec:conclusiones} presenta las conclusiones y recomendaciones de política pública.

% ============================================================
\section{Marco Teórico y Revisión de Literatura}
\label{sec:marco}
% ============================================================

\subsection{Informalidad laboral: definición y medición}

La informalidad laboral se define, en su acepción más operacional, como toda relación
de trabajo que no está sujeta a la legislación nacional, no cumple con el pago de
impuestos y carece de cobertura de protección social. Puede presentarse a nivel de persona
(trabajador que no cotiza) o a nivel de empresa (unidad productiva no registrada).

Para este estudio se adopta la definición operativa del DANE: un trabajador es
\textbf{informal} (INFORMAL = 1) si en la variable P6920 de la GEIH declara no cotizar
a salud mediante su actividad laboral. Los pensionados (P6920 = 3) se excluyen porque
su cobertura previsional proviene de la pensión, no del trabajo actual, lo que introduciría
sesgo en el modelo.

\begin{equation}
  \text{INFORMAL}_i =
  \begin{cases}
    0 & \text{si } P6920_i = 1 \text{ (cotizante)}\\
    1 & \text{si } P6920_i = 2 \text{ (no cotizante)}\\
    \text{NaN} & \text{si } P6920_i = 3 \text{ (pensionado — excluido)}
  \end{cases}
\end{equation}

\subsection{Contexto colombiano}

Colombia muestra una tasa de informalidad ponderada (ajustada por factores de expansión
DANE) del \textbf{56{,}0\%} para la población ocupada general. Entre los trabajadores
independientes (cuenta propia) esta tasa asciende al \textbf{84{,}5\%}: ser trabajador
por cuenta propia es, en el contexto colombiano, casi sinónimo de ser informal.

La Ley 1429 de 2010 intentó incentivar la formalización mediante beneficios tributarios
y descuentos en parafiscales para pequeñas empresas. Sin embargo, las evaluaciones de
impacto muestran efectos heterogéneos: promovió la formalización en algunos segmentos
pero generó preferencia por mano de obra poco calificada en otros, aumentando la
informalidad en ciertos subgrupos juveniles \cite{arizacedano2017}.

\subsection{Machine Learning en predicción de informalidad}

La literatura reciente ha utilizado modelos supervisados para estimar probabilidades
individuales de informalidad. Los modelos basados en árboles de decisión (Random Forest,
Gradient Boosting) superan consistentemente a los modelos lineales en datos tabulares
estructurados \cite{grinsztajn2022}. En particular, LightGBM muestra ventajas de velocidad
y rendimiento en conjuntos de datos de tamaño mediano (100K--1M filas) característicos
de encuestas nacionales.

Para la interpretabilidad se utiliza SHAP (SHapley Additive exPlanations), que asigna
a cada variable una contribución marginal a cada predicción individual, satisfaciendo las
propiedades de Shapley: eficiencia, simetría, nulidad y linealidad \cite{lundberg2017}.

\subsection{Metodología: CRISP-DM}

El proyecto sigue la metodología CRISP-DM (Cross-Industry Standard Process for Data
Mining) \cite{crisp2000}, que organiza el trabajo en seis fases cíclicas:
entendimiento del negocio, entendimiento de los datos, preparación de los datos,
modelado, evaluación y despliegue. Esta metodología fue elegida por su orientación
al problema de negocio (focalización de política pública) y por su trazabilidad:
cada hallazgo del EDA justifica una decisión en la preparación, y cada métrica del
modelo se conecta con el objetivo de focalización.

% ============================================================
\section{Desarrollo Metodológico — Machine Learning}
\label{sec:ml}
% ============================================================

\subsection{Entendimiento del problema y evolución de la pregunta de investigación}

\subsubsection{Pregunta inicial y refinamiento}

La pregunta de investigación inicial era amplia:
\textit{``¿Cuáles son las características socioeconómicas y laborales que mejor predicen la
informalidad laboral en Colombia?''}

Durante el análisis exploratorio se identificó que el grupo de trabajadores
\textbf{independientes (cuenta propia, P6430 = 4)} concentra el problema de una forma
cualitativamente distinta: su tasa de informalidad es 84{,}5\% (vs.\ 56\% del promedio),
y las variables que predicen su informalidad son distintas a las del trabajador asalariado
(para el asalariado, el tipo de contrato es el predictor dominante; para el independiente,
el tamaño del establecimiento y la educación son más relevantes).

Esto llevó a refinar la pregunta hacia:
\textit{``¿Por qué los trabajadores independientes en Colombia permanecen informales,
y qué condiciones deben cambiar para que la formalización sea una opción viable para ellos?''}

Esta decisión implicó filtrar el dataset a P6430 = 4 (151.136 registros), un cambio
que mejora la especificidad del modelo y su utilidad para el diseño de políticas
orientadas al segmento de mayor vulnerabilidad.

\subsubsection{Variables excluidas por diseño}

\begin{decision}
La variable P6450 (tipo de contrato: verbal / escrito) fue excluida del modelo
para trabajadores independientes. Un trabajador por cuenta propia que declara
``contrato verbal'' está describiendo una relación atípica para su categoría:
el contrato es, por definición, proxy de la informalidad misma. Incluirlo
generaría un modelo trivialmente correcto pero vacío de contenido causal.
La misma lógica aplica a CONTRATO\_VERBAL (variable derivada de P6450).
\end{decision}

La Tabla~\ref{tab:excluidas} resume las variables excluidas y su justificación.

\begin{table}[H]
\centering
\caption{Variables excluidas del modelo y justificación}
\label{tab:excluidas}
\begin{tabularx}{\textwidth}{lX}
\toprule
\textbf{Variable} & \textbf{Razón de exclusión} \\
\midrule
INGLABO (ingreso laboral) & Data leakage: consecuencia de la informalidad, no causa \\
P6460 (naturaleza contrato) & 60{,}78\% de nulos — imputación estadísticamente indefendible \\
P6920 & Fuente de la variable objetivo \\
FEX\_C18 & Factor de expansión muestral, no predictor \\
P6450 / CONTRATO\_VERBAL & Proxy del target para trabajadores independientes \\
P6430 / CUENTA\_PROPIA & Variable de filtro; todas las obs. tienen P6430 = 4 \\
\bottomrule
\end{tabularx}
\end{table}

\subsection{Análisis Exploratorio de Datos (EDA)}

\subsubsection{Fuente y carga de datos}

La \textbf{Gran Encuesta Integrada de Hogares 2024} (GEIH) del DANE cubre los 12 meses
del año en dos módulos CSV por mes: \textit{Características Generales} y \textit{Ocupados}.
En total se procesan 24 archivos CSV con codificación \texttt{latin-1} y separador punto
y coma. La llave de unión entre módulos es la tripleta
\texttt{DIRECTORIO + SECUENCIA\_P + ORDEN}, que identifica unívocamente a cada persona.

\begin{table}[H]
\centering
\caption{Volúmenes de datos a través del pipeline}
\label{tab:volumenes}
\begin{tabular}{lrr}
\toprule
\textbf{Etapa} & \textbf{Registros} & \textbf{Observación} \\
\midrule
Módulo Características Generales & 817.550 & Total encuestados \\
Módulo Ocupados & 358.029 & Solo población ocupada \\
Tras inner join & 358.029 & Solo los presentes en ambos módulos \\
Excluir pensionados (P6920 = 3) & 351.721 & 1{,}8\% excluidos \\
Filtro cuenta propia (P6430 = 4) & 151.136 & Población objetivo \\
Train (80\%, estratificado) & 120.909 & Para ajuste del modelo \\
Test (20\%, estratificado) & 30.227 & Solo para evaluación final \\
\bottomrule
\end{tabular}
\end{table}

\subsubsection{Tasa de informalidad ponderada}

La GEIH no es un muestreo uniforme: cada persona tiene un factor de expansión
\texttt{FEX\_C18} que indica cuántas personas de la población representa. La tasa
ponderada correcta es:

\begin{equation}
  \hat{\tau} = \frac{\sum_i \text{INFORMAL}_i \cdot \text{FEX\_C18}_i}
                    {\sum_i \text{FEX\_C18}_i} = 56{,}0\%
  \quad \text{(todos los ocupados)}
\end{equation}

Para trabajadores por cuenta propia: $\hat{\tau} = 84{,}5\%$. El desbalanceo de clases
es mayor en este subgrupo (razón 5{,}5:1), pero sigue siendo manejable con el parámetro
\texttt{is\_unbalance=True} de LightGBM.

\subsubsection{Hallazgos principales del EDA}

Los siguientes hallazgos del EDA justifican directamente las decisiones de ingeniería
de características:

\textbf{Educación (P3042):} La tasa de informalidad cae de $\approx$90\% para trabajadores
sin educación formal a $\approx$15\% para doctores. La relación es negativa y casi
monotónica, pero los intervalos entre niveles no son iguales (pasar de ``ninguno'' a
``preescolar'' no equivale a pasar de ``universitaria'' a ``maestría''). Esto justifica
la creación de \textit{ANIOS\_EDU} como variable continua de años de escolaridad acumulados.

\textbf{Edad (P6040):} La curva de informalidad por edad tiene forma de U: mínimo en el
tramo 25--34 años (inserción al mercado formal) y repunte fuerte después de los 64
(personas mayores que trabajan informalmente después de la edad de retiro). Esta
no-linealidad justifica mantener tanto la edad continua (para el modelo de árboles)
como el grupo etario categórico \textit{EDAD\_GRUPO} (para el OHE).

\textbf{Tamaño del establecimiento (P3069):} Establecimientos con 1--10 personas tienen
tasas de informalidad de 75--95\%. La caída es drástica a partir de 20 personas. Esto
justifica la variable binaria \textit{MICROEMPRESA} ($P3069 \in \{1, 2, 3\}$).

\begin{hallazgo}
\textbf{Interacción Microempresa $\times$ Tipo de contrato.} El análisis bivariado revela
una interacción crucial: en microempresas, el contrato verbal se asocia con $\approx$95\%
de informalidad, pero el contrato escrito la reduce a $\approx$60\% (-35 puntos
porcentuales). En grandes empresas, la misma transición reduce de $\approx$35\% a $\approx$8\%
(-27 pp). El efecto protector del contrato escrito es \textbf{mayor} donde más informalidad
existe (las microempresas). Esta es la recomendación de política pública más accionable
que emerge del análisis.
\end{hallazgo}

\subsection{Preparación de los datos e ingeniería de características}

\subsubsection{Variables derivadas}

La Tabla~\ref{tab:features_eng} muestra las siete variables derivadas y su justificación
en el EDA.

\begin{table}[H]
\centering
\caption{Ingeniería de características — variables derivadas}
\label{tab:features_eng}
\begin{tabularx}{\textwidth}{llX}
\toprule
\textbf{Variable} & \textbf{Derivación} & \textbf{Justificación EDA} \\
\midrule
ANIOS\_EDU & Mapa P3042$\to$años + P3042S1 & Relación negativa continua con informalidad \\
EDAD\_GRUPO & \texttt{pd.cut}(P6040, 6 tramos) & Curva en U no capturada por edad continua sola \\
MICROEMPRESA & P3069 $\in$ \{1,2,3\} $\to$ 0/1 & Tasas 75--95\% en establec. de 1--10 personas \\
SUBEMPLEADO & P6800 $<$ 32 $\to$ 0/1 & Menos de 32h/sem = subempleo (80\% jornada legal) \\
PLURIEMPLEO & P7040: \{1$\to$1, 2$\to$0\} & Recodificación de binario GEIH \\
\bottomrule
\end{tabularx}
\end{table}

\subsubsection{Pipeline de preprocesamiento (scikit-learn)}

El corazón de la preparación es un \texttt{ColumnTransformer} que aplica tratamientos
diferenciados según el tipo estadístico de cada variable:

\begin{lstlisting}[caption={Pipeline ColumnTransformer — extracto de 3.Preparacion.ipynb}]
from sklearn.pipeline import Pipeline
from sklearn.compose import ColumnTransformer
from sklearn.preprocessing import StandardScaler, OneHotEncoder
from sklearn.impute import SimpleImputer
from category_encoders import TargetEncoder

FEATURES = [
    "P3271", "P6040", "EDAD_GRUPO", "ANIOS_EDU", "P6070",
    "CLASE", "DPTO", "P6800", "P3069", "RAMA2D_R4",
    "MICROEMPRESA", "SUBEMPLEADO", "PLURIEMPLEO"
]

num_pipe   = Pipeline([("imp", SimpleImputer(strategy="median")),
                        ("scl", StandardScaler())])
bin_pipe   = Pipeline([("imp", SimpleImputer(strategy="most_frequent"))])
ohe_pipe   = Pipeline([("imp", SimpleImputer(strategy="most_frequent")),
                        ("ohe", OneHotEncoder(handle_unknown="ignore",
                                              sparse_output=False))])
tgt_pipe   = Pipeline([("imp", SimpleImputer(strategy="most_frequent")),
                        ("tgt", TargetEncoder(smoothing=1.0))])

preprocessor = ColumnTransformer([
    ("num", num_pipe, COLS_NUM),
    ("bin", bin_pipe, COLS_BIN),
    ("ohe", ohe_pipe, COLS_OHE),
    ("tgt", tgt_pipe, COLS_TGT),
], remainder="drop")
\end{lstlisting}

\begin{decision}
\textbf{¿Por qué TargetEncoder para RAMA2D\_R4 y DPTO?} Ambas variables tienen alta
cardinalidad (40 ramas CIIU, 33 departamentos). OneHotEncoding generaría hasta 73
columnas adicionales, muchas con menos de 100 registros. TargetEncoder las convierte
en un único número (probabilidad media de informalidad para esa categoría) con suavizado
$\lambda = 1{,}0$ que estabiliza estimaciones para categorías pequeñas (ej. departamento
Vaupés con pocos registros).
\end{decision}

\begin{decision}
\textbf{¿Por qué dividir train/test \textit{antes} de ajustar el pipeline?} Si el
preprocesador se ajustara sobre todo el dataset y luego se dividiera, el conjunto de
prueba habría ``informado'' las medias y escalas del preprocesador, generando fuga de
información (\textit{data leakage}). La secuencia correcta es:
(1) \texttt{train\_test\_split(stratify=y)}, (2) \texttt{preprocessor.fit(X\_train)},
(3) \texttt{preprocessor.transform(X\_test)}.
\end{decision}

\subsection{Selección y entrenamiento de modelos}

\subsubsection{Criterio de evaluación: F1-Score}

Se selecciona el F1-Score como métrica principal por dos razones:

\begin{enumerate}[label=\arabic*.]
  \item Hay desbalanceo de clases (84{,}5\% informal vs.\ 15{,}5\% formal en cuenta
  propia): la exactitud (\textit{accuracy}) sería engañosa (un clasificador trivial
  ``siempre informal'' obtendría 84{,}5\%).
  \item Tanto los falsos positivos (clasificar formal a un informal) como los falsos
  negativos (perder a un informal real) tienen costos de política: el primero malgasta
  recursos en quien no los necesita; el segundo excluye a quien más los necesita.
  F1 = media armónica entre precisión y recall equilibra ambos errores.
\end{enumerate}

\subsubsection{Los cuatro modelos}

Se entrenaron cuatro modelos siguiendo el principio CRISP-DM de comparar alternativas
antes de seleccionar el campeón:

\begin{table}[H]
\centering
\caption{Comparativa de modelos — conjunto de prueba (n = 70.345)}
\label{tab:modelos}
\begin{tabular}{lcccc}
\toprule
\textbf{Modelo} & \textbf{F1} & \textbf{AUC-ROC} & \textbf{Precisión} & \textbf{Recall} \\
\midrule
\rowcolor{azulDane!10}
\textbf{LightGBM (tuned)} & \textbf{0{,}9576} & \textbf{0{,}9853} & 0{,}9399 & 0{,}9760 \\
Random Forest & 0{,}9553 & 0{,}9818 & 0{,}9378 & 0{,}9734 \\
XGBoost & 0{,}9546 & 0{,}9825 & 0{,}9423 & 0{,}9673 \\
Regresión Logística & 0{,}9507 & 0{,}9738 & 0{,}9347 & 0{,}9673 \\
\bottomrule
\end{tabular}
\end{table}

La Regresión Logística actúa como \textit{baseline}: establece el piso de rendimiento
de un modelo lineal. Random Forest y XGBoost son los competidores de ensemble.
LightGBM es el candidato principal por su velocidad con datos medianos-grandes y su
historial en competencias de clasificación tabular.

\begin{decision}
\textbf{¿Por qué LightGBM sobre XGBoost?} LightGBM crece las hojas del árbol con mayor
ganancia (\textit{leaf-wise}) en lugar de crecer nivel a nivel (\textit{level-wise} de
XGBoost), permitiendo mayor profundidad efectiva con menos iteraciones. En este dataset
LightGBM entrenó en $\approx$40 s vs.\ $\approx$3 min de XGBoost con configuraciones
comparables. El F1 fue 0{,}0030 mayor (diferencia pequeña pero consistente en las cuatro
métricas). Para datos tabulares estructurados, la literatura confirma que los modelos
de árboles superan a redes neuronales \cite{grinsztajn2022}.
\end{decision}

\subsubsection{Estrategia de entrenamiento LightGBM en tres pasos}

\textbf{Paso 1 — Early stopping.} Se separa el 15\% del conjunto de entrenamiento como
validación temporal. Se entrena con hasta 5.000 árboles pero se detiene si 100 iteraciones
consecutivas no mejoran el \texttt{binary\_logloss} en validación. Resultado: iteración
óptima = 858. Esto evita el sobreajuste sin búsqueda exhaustiva de hiperparámetros.

\textbf{Paso 2 — Threshold tuning.} El umbral por defecto (0{,}5) convierte probabilidades
en clases. Sin embargo, el F1 óptimo no necesariamente ocurre en 0{,}5. Se busca el umbral
que maximiza el F1 sobre el conjunto de validación en el rango [0{,}30; 0{,}70].
\textbf{Resultado: umbral óptimo = 0{,}41.} El modelo es ligeramente más propenso a
clasificar como informal, lo que es coherente con la alta prevalencia de informalidad
(84{,}5\%) en el grupo objetivo.

\begin{equation}
  t^* = \arg\max_{t \in [0{,}30, 0{,}70]}
        F_1\!\left(\hat{y}(t),\ y_\text{val}\right)
  \quad \Rightarrow \quad t^* = 0{,}41
\end{equation}

\textbf{Paso 3 — Reentrenamiento final.} Con el número óptimo de iteraciones (858),
se reentrena el modelo sobre el 100\% del conjunto de entrenamiento. Esto aprovecha
todos los datos disponibles para el modelo final sin riesgo de sobreajuste (el número
de iteraciones ya fue calibrado con \textit{early stopping}).

\subsubsection{Hiperparámetros clave}

\begin{lstlisting}[caption={Hiperparámetros LightGBM campeón}]
PARAMS_LGBM = dict(
    learning_rate    = 0.02,    # Bajo: mas arboles, mejor generalizacion
    num_leaves       = 127,     # Complejidad moderada (2^7 - 1)
    max_depth        = -1,      # Sin limite de profundidad
    min_child_samples = 40,     # Minimo de obs. por hoja
    reg_alpha        = 0.1,     # Regularizacion L1
    reg_lambda       = 1.0,     # Regularizacion L2
    colsample_bytree = 0.7,     # Submuestreo de features por arbol
    colsample_bynode = 0.7,     # Submuestreo por nodo (reduce correlacion)
    subsample        = 0.8,     # Submuestreo de observaciones
    subsample_freq   = 1,
    is_unbalance     = True,    # Pesos inversos a frecuencia de clase
    n_jobs           = -1,
    random_state     = 42,
    verbose          = -1,
)
\end{lstlisting}

\subsection{Evaluación del modelo}

\subsubsection{Métricas finales — conjunto de prueba}

\begin{table}[H]
\centering
\caption{Métricas del modelo campeón (LightGBM) en test set}
\label{tab:metricas_final}
\begin{tabular}{lcc}
\toprule
\textbf{Métrica} & \textbf{Valor} & \textbf{Interpretación} \\
\midrule
F1-Score & \textbf{0{,}9576} & Equilibrio precisión-recall \\
AUC-ROC & \textbf{0{,}9853} & Calidad del ranking de probabilidades \\
Precisión & 0{,}9399 & 93{,}99\% de los clasificados informales lo son \\
Recall & 0{,}9760 & Detecta el 97{,}60\% de los informales reales \\
Umbral de decisión & 0{,}41 & Calibrado para maximizar F1 \\
Iteraciones (árboles) & 858 & Determinado por early stopping \\
\bottomrule
\end{tabular}
\end{table}

Un AUC-ROC de 0{,}9853 significa que, si se toma al azar un informal y un formal,
el modelo asigna mayor probabilidad de informalidad al informal en el \textbf{98{,}53\%}
de las veces — una medida robusta independiente del umbral de decisión.

\subsubsection{Análisis SHAP — Interpretabilidad}

SHAP (TreeExplainer) calcula valores exactos para modelos de árboles en
$O(\text{features} \times \text{profundidad})$, siendo $>100\times$ más rápido que
KernelExplainer (Monte Carlo) sin pérdida de exactitud. La muestra utilizada es de
3.000 observaciones del test set, suficiente para una estimación estable de la
importancia media.

\begin{table}[H]
\centering
\caption{Top 10 variables — importancia SHAP media $|\phi_j|$}
\label{tab:shap}
\begin{tabular}{rlcc}
\toprule
\textbf{Rank} & \textbf{Variable} & \textbf{SHAP medio} & \textbf{Interpretación} \\
\midrule
1 & Microempresa ($\leq$10 personas) & 1{,}397 & Factor más determinante de informalidad \\
2 & Años de educación & 0{,}576 & Más educación = menos informalidad \\
3 & P3069\_10 (No sabe tamaño) & 0{,}546 & Desconocer tamaño = perfil informal \\
4 & Departamento (enc. objetivo) & 0{,}404 & Geografía importa (Vaupés 84{,}3\% vs.\ Bogotá 26\%) \\
5 & P3069\_1 (1 sola persona) & 0{,}361 & Autoempleo unipersonal = informal \\
6 & Edad & 0{,}329 & Curva en U: jóvenes y mayores más vulnerables \\
7 & Horas / semana & 0{,}326 & Pocas horas $\to$ subempleo informal \\
8 & Rama de actividad (enc. obj.) & 0{,}302 & Construcción, hogar, agric. = informales \\
9 & Estado civil & 0{,}221 & Efecto social / responsabilidad familiar \\
10 & Zona (urbano/rural) & 0{,}198 & Rural $\to$ mayor informalidad \\
\bottomrule
\end{tabular}
\end{table}

\begin{hallazgo}
Microempresa es el predictor más importante con una ventaja de 2{,}4× sobre el
segundo lugar (educación). Esto indica que el tamaño del establecimiento donde trabaja
el independiente concentra la mayor información para predecir su condición de
informalidad — por encima de su perfil demográfico o su ubicación geográfica.
\end{hallazgo}

\subsection{Análisis, implicaciones y conclusiones del módulo ML}

El modelo confirma que la informalidad en trabajadores independientes no es aleatoria
sino el resultado predecible de condiciones estructurales:

\begin{enumerate}[label=\arabic*.]
  \item \textbf{Tamaño del establecimiento} ($\phi = 1{,}397$): trabajar en una unidad
  productiva de menos de 10 personas es el mayor predictor de informalidad, lo que
  refleja la baja capacidad administrativa de la microempresa para cumplir obligaciones
  parafiscales.
  \item \textbf{Educación} ($\phi = 0{,}576$): cada año adicional de educación reduce
  la probabilidad de informalidad. La relación es causal en el sentido de que la educación
  expande el acceso a mercados laborales formales y la capacidad de negociación del trabajador.
  \item \textbf{Geografía} ($\phi = 0{,}404$): la informalidad en Vaupés (84{,}3\%) es
  $3{,}2\times$ mayor que en Bogotá ($\approx$26\%), lo que refleja diferencias en densidad
  empresarial, acceso a servicios del Estado y estructura económica departamental.
\end{enumerate}

% ============================================================
\section{Tecnología e Ingeniería de Datos}
\label{sec:tecnologia}
% ============================================================

\subsection{Arquitectura de referencia}

El proyecto sigue una arquitectura \textbf{batch analítica local}, apropiada para el
volumen de datos (6{,}2 MB en Parquet) y el objetivo académico-investigativo.
En un escenario de implementación real a escala nacional, la arquitectura escalaría
hacia una solución de nube híbrida.

\begin{figure}[H]
\centering
\fbox{\parbox{0.90\textwidth}{\centering\small
\texttt{DANE GEIH 2024 (24 CSV, 180 MB)}\\
$\downarrow$ Ingesta batch (Python pandas)\\
\texttt{parquet/geih\_2024\_crudo.parquet} + \texttt{geih\_2024.duckdb} (6{,}2 MB)\\
$\downarrow$ EDA exploratorio\\
\texttt{2.EDA.ipynb} (sin artefactos de salida — solo análisis)\\
$\downarrow$ Feature engineering + Pipeline ML\\
\texttt{parquet/train.parquet} · \texttt{test.parquet} · \texttt{preprocessor.joblib}\\
$\downarrow$ Modelado y evaluación\\
\texttt{outputs/champion\_geih.pkl} · \texttt{meta.json} · \texttt{metrics.csv} · \texttt{*.png}\\
$\downarrow$ Visualizaciones complementarias\\
\texttt{outputs/viz\_01…viz\_06.html}\\
$\downarrow$ Aplicación web\\
\texttt{app.py} (Streamlit) $\to$ Streamlit Community Cloud (enlace público)
}}
\caption{Pipeline de datos — arquitectura de referencia del proyecto}
\label{fig:pipeline}
\end{figure}

\subsection{Origen y naturaleza de los datos}

\begin{table}[H]
\centering
\caption{Descripción de los datasets utilizados}
\label{tab:datasets}
\begin{tabularx}{\textwidth}{lllX}
\toprule
\textbf{Dataset} & \textbf{Tipo} & \textbf{Modo ingesta} & \textbf{Descripción} \\
\midrule
GEIH 2024 — Caract.\ Generales & Estructurado & Batch mensual & 12 CSV: sexo, edad, educación, estado civil, zona, departamento \\
GEIH 2024 — Ocupados & Estructurado & Batch mensual & 12 CSV: posición ocupacional, horas, tamaño empresa, rama, contrato \\
\bottomrule
\end{tabularx}
\end{table}

Los datos son completamente \textbf{estructurados}, con separador \texttt{;} y
codificación \texttt{latin-1} (típico de los archivos DANE). No hay datos en streaming,
transaccionales ni no estructurados en este proyecto — lo cual es coherente con el origen
estadístico-muestral de la GEIH.

\subsection{Ingesta de datos (modo batch)}

La ingesta se realiza en el notebook \texttt{1.Ingesta.ipynb} con una función
\texttt{cargar\_modulo()} que procesa los 12 archivos mensuales de cada módulo en lote:

\begin{lstlisting}[caption={Función de ingesta batch — 1.Ingesta.ipynb}]
def cargar_modulo(carpeta, columnas_interes):
    """Carga 12 archivos CSV mensuales, filtra columnas y concatena."""
    dfs = []
    for archivo in sorted(Path(carpeta).glob("*.CSV")):
        df = pd.read_csv(archivo, sep=";", encoding="latin-1",
                         usecols=columnas_interes, low_memory=False)
        dfs.append(df)
    return pd.concat(dfs, ignore_index=True)
\end{lstlisting}

Solo se cargan las 19 columnas relevantes identificadas en el diseño, lo que reduce
el footprint de memoria en $6\times$--$10\times$ respecto a cargar todos los CSV completos.

\subsection{Almacenamiento: Parquet y DuckDB}

\subsubsection{¿Por qué Parquet?}

El formato Parquet (Apache, columnar, comprimido con Snappy) ofrece:
\begin{itemize}[noitemsep]
  \item \textbf{Compresión}: 180 MB en CSV $\to$ 6{,}2 MB en Parquet (ratio 29:1).
  \item \textbf{Acceso columnar}: leer solo las columnas necesarias sin deserializar
  todo el archivo — crítico para datasets anchos.
  \item \textbf{Tipado estricto}: preserva dtypes (int32, float32) sin conversiones
  implícitas de CSV.
  \item \textbf{Interoperabilidad}: compatible con Pandas, DuckDB, Spark, y cualquier
  motor analítico moderno.
\end{itemize}

\subsubsection{¿Por qué DuckDB?}

DuckDB es un motor SQL OLAP \textit{in-process} (sin servidor). Se ejecuta directamente
en el proceso Python, soporta SQL analítico completo sobre archivos Parquet, y es
el equivalente \textit{embedded} de un sistema como BigQuery o Redshift para análisis
ad-hoc locales. Su uso cumple el criterio SI7006 de ``motores SQL/NoSQL de acuerdo
con la problemática'':

\begin{lstlisting}[caption={Consulta analítica en DuckDB sobre el Parquet}]
import duckdb
conn = duckdb.connect("geih_2024.duckdb")
conn.execute("CREATE TABLE IF NOT EXISTS geih AS "
             "SELECT * FROM 'parquet/geih_2024_crudo.parquet'")
# Ejemplo: tasa de informalidad por departamento
tasa_dpto = conn.execute("""
    SELECT DPTO,
           SUM(INFORMAL * FEX_C18) / SUM(FEX_C18) AS tasa_inf
    FROM geih
    WHERE INFORMAL IS NOT NULL
    GROUP BY DPTO ORDER BY tasa_inf DESC
""").df()
\end{lstlisting}

\subsection{Ambiente de procesamiento}

\begin{table}[H]
\centering
\caption{Stack tecnológico del proyecto}
\label{tab:stack}
\begin{tabular}{lll}
\toprule
\textbf{Componente} & \textbf{Herramienta} & \textbf{Versión} \\
\midrule
Lenguaje & Python & 3{,}11{,}9 \\
Entorno virtual & venv (local) & — \\
Análisis y manipulación & Pandas & 2{,}2{,}x \\
Cómputo numérico & NumPy & 2{,}4{,}x \\
ML pipeline & scikit-learn & 1{,}4{,}x \\
Modelo campeón & LightGBM & 4{,}6{,}0 \\
Interpretabilidad & SHAP & 0{,}51{,}0 \\
Almacenamiento columnar & PyArrow / Parquet & 15{,}x \\
Motor SQL & DuckDB & 0{,}10{,}x \\
Visualización & Plotly & 5{,}18{,}x \\
Dashboard & Streamlit & 1{,}58{,}0 \\
Recomendaciones IA & Groq (Llama 3.3 70B) & API v1 \\
\bottomrule
\end{tabular}
\end{table}

\subsection{Persistencia y serialización del modelo}

El modelo campeón se persiste en dos formatos complementarios:

\begin{lstlisting}[caption={Serialización del modelo y metadatos — 4.Modelamiento.ipynb}]
import pickle, json

# Modelo
with open("outputs/champion_geih.pkl", "wb") as f:
    pickle.dump(model_final, f)    # 11.48 MB

# Metadatos: metricas + umbral + lista de features
meta = {
    "model": "LightGBM",
    "threshold": 0.41,
    "test_f1": 0.9576,
    "test_auc": 0.9853,
    "test_precision": 0.9399,
    "test_recall": 0.9760,
    "n_train": 120909,
    "n_test": 30227,
    "n_features": len(feature_names),
    "feature_names": feature_names,
}
with open("outputs/champion_geih_meta.json", "w") as f:
    json.dump(meta, f, indent=2)
\end{lstlisting}

El archivo JSON de metadatos desacopla las métricas del objeto modelo: la interfaz
(app.py) puede mostrar F1, AUC y umbral sin cargar el modelo de 11 MB.

\subsection{Despliegue del modelo — escenario real vs.\ académico}

\textbf{Escenario académico (implementado):} El modelo se despliega como aplicación
Streamlit en Streamlit Community Cloud. El código y los artefactos residen en un
repositorio GitHub público. Streamlit Cloud clona el repositorio, instala las
dependencias de \texttt{requirements.txt}, y ejecuta \texttt{app.py}. El enlace
resultante es accesible públicamente sin autenticación.

\textbf{Escenario hipotético de implementación real (Ministerio del Trabajo):}
En un contexto de adopción institucional, la arquitectura escalaría así:
\begin{itemize}[noitemsep]
  \item \textbf{Datos}: ingesta mensual automatizada desde el FTP del DANE
  hacia un Data Lake en S3 o GCS, con particionamiento por mes y módulo.
  \item \textbf{Procesamiento}: pipeline ETL en Apache Spark sobre EMR/Dataproc
  para el join mensual de los módulos ($\approx$70 millones de registros/año).
  \item \textbf{Modelo}: servido como microservicio REST (FastAPI + Docker),
  con reentrenamiento automático mensual al recibir nuevos datos GEIH.
  \item \textbf{Dashboard}: Power BI o Streamlit Enterprise detrás de
  autenticación LDAP del Ministerio, con control de acceso por rol.
  \item \textbf{MLOps}: MLflow para tracking de experimentos y comparación
  automática de modelos nuevos vs.\ modelo en producción.
\end{itemize}

\subsection{Optimización de rendimiento en el dashboard}

El dashboard utiliza el sistema de caché de Streamlit para evitar recargas costosas:

\begin{lstlisting}[caption={Estrategia de caché en app.py}]
@st.cache_resource          # Objeto Python costoso en memoria
def load_model():
    with open(MODEL_PATH, "rb") as f:
        return pickle.load(f)    # 11.48 MB, ~2 seg carga

@st.cache_data              # DataFrame serializable
def load_data():
    return pd.read_parquet(DATA_PATH)   # ~500 ms carga

@st.cache_data(ttl=3600)    # SHAP: costoso de computar (~10-30 seg)
def compute_shap_beeswarm():
    model = load_model()
    # ... calculo TreeExplainer sobre 800 muestras ...
    return dict_serializable    # Resultados cacheados por 1 hora
\end{lstlisting}

\texttt{cache\_resource} mantiene el modelo en memoria entre recargas de la página.
\texttt{cache\_data} serializa y cachea DataFrames. Sin este mecanismo, cada
interacción del usuario recargaría el modelo (2 seg) y el dataset (500 ms), haciendo
la experiencia inaceptablemente lenta.

% ============================================================
\section{Visualización y Comunicación de Datos}
\label{sec:viz}
% ============================================================

\subsection{Requerimientos de comunicación}

El diseño del dashboard parte de una pregunta concreta: \textit{¿qué necesita
ver un funcionario del Ministerio del Trabajo o del SENA para tomar una decisión
de focalización?} A partir de esta pregunta se identificaron tres requerimientos
funcionales de comunicación:

\begin{enumerate}[label=\arabic*.]
  \item \textbf{Magnitud y distribución del problema.}
  El usuario debe poder responder en menos de 30 segundos: ¿qué tan grave es la
  informalidad entre los independientes?, ¿dónde está concentrada geográficamente?
  Esto demanda visualizaciones de resumen con cifras absolutas y relativas,
  ordenadas para que los valores extremos sean inmediatamente visibles.

  \item \textbf{Factores determinantes — comprensibles para no técnicos.}
  La política pública requiere entender \textit{por qué} un trabajador es informal,
  no solo \textit{que} lo es. Las visualizaciones deben traducir los resultados del
  modelo (valores SHAP, importancia de variables) a lenguaje de negocio: qué
  características elevan el riesgo, con qué magnitud, y qué combinaciones son
  más críticas. Los nombres de variables deben ser comprensibles sin conocer el
  diccionario GEIH.

  \item \textbf{Capacidad de exploración individual.}
  El caso de uso operativo más valioso es el del inspector o gestor que tiene
  frente a sí un perfil de trabajador y necesita estimar su riesgo de informalidad.
  El dashboard debe permitir ingresar características individuales y obtener una
  probabilidad, un veredicto y recomendaciones de acción concretas —
  sin requerir conocimiento estadístico ni de Machine Learning.
\end{enumerate}

\subsection{Elección de herramientas}

\subsubsection{¿Por qué Streamlit y no Dash o Power BI?}

Streamlit fue elegido sobre Dash y Power BI por las siguientes razones:

\begin{itemize}[noitemsep]
  \item \textbf{Integración nativa con Python}: el modelo, el preprocesador, y los datos
  están en Python. Streamlit ejecuta código Python directamente sin capa de traducción.
  \item \textbf{Despliegue gratuito}: Streamlit Community Cloud permite despliegue
  desde GitHub sin costo, con URL pública persistente.
  \item \textbf{Velocidad de prototipado}: una visualización que en Dash requeriría
  definir callbacks y layouts HTML, en Streamlit se escribe en 5 líneas de Python.
  \item \textbf{Curva de aprendizaje}: el equipo ya dominaba Python; aprender Dash
  (React/JavaScript subyacente) o Power BI (lenguaje M, DAX) habría sido una inversión
  desproporcionada para el alcance del proyecto.
\end{itemize}

\subsubsection{¿Por qué Plotly y no Matplotlib/Seaborn?}

\begin{itemize}[noitemsep]
  \item \textbf{Interactividad}: Plotly genera gráficas con zoom, hover, descarga PNG,
  y selección de trazas — todo sin JavaScript adicional.
  \item \textbf{Integración con Streamlit}: \texttt{st.plotly\_chart()} renderiza objetos
  Plotly nativamente. Matplotlib genera imágenes estáticas (PNG) que perderían toda la
  interactividad.
  \item \textbf{Calidad visual}: Plotly produce gráficas de calidad publicación sin
  configuración adicional; Seaborn requiere mucho ajuste para alcanzar el mismo nivel.
\end{itemize}

\subsection{Diseño y paleta de colores}

\subsubsection{Paleta base: escala de grises con acento semántico}

El dashboard usa una paleta de alto contraste sobre fondo blanco:

\begin{itemize}[noitemsep]
  \item \textbf{Negro} (\texttt{\#222222}): texto principal, títulos, ejes.
  \item \textbf{Gris claro} (\texttt{\#E0E0E0}): barras de referencia, fondos de gráfica.
  \item \textbf{Colores semánticos de riesgo}: solo para indicar nivel de informalidad.
    \begin{itemize}[noitemsep, label=$-$]
      \item Rojo oscuro (\texttt{\#C0392B}): tasa $\geq$ promedio + 8 pp (zona crítica).
      \item Naranja (\texttt{\#E67E22}): tasa $\geq$ promedio.
      \item Amarillo (\texttt{\#F1C40F}): tasa moderada.
      \item Verde (\texttt{\#27AE60}): tasa baja (zona segura).
    \end{itemize}
\end{itemize}

\begin{decision}
\textbf{Justificación de la paleta.} El uso de escala de grises como base y colores
semánticos solo para riesgo sigue el principio de \textit{data-ink ratio} de Tufte:
el color solo se añade cuando agrega información, no como decoración. Un usuario que
ve una barra roja inmediatamente sabe que representa alto riesgo, sin necesidad de
leer la leyenda. Los colores semánticos refuerzan la narrativa (``zona crítica'' en rojo,
``zona segura'' en verde) y son accesibles para personas con daltonismo parcial
(la secuencia rojo-naranja-amarillo-verde es distinguible incluso en deuteranopía
cuando se combina con variaciones de luminosidad).
\end{decision}

\subsection{Análisis de las visualizaciones por tab}

\subsubsection{Tab 1 — Arquitectura \& Pipeline}

Presenta el flujo de datos como un diagrama de árbol ASCII embebido en Markdown,
complementado con una tabla de stack tecnológico. La visualización es deliberadamente
estática y esquemática: su audiencia es el evaluador de arquitectura, quien necesita
leer texto estructurado, no interactuar con gráficas. Un diagrama de nodos-aristas
animado sería visualmente más atractivo pero reduciría la densidad de información.

\subsubsection{Tab 2 — Dashboard ¿Por qué permanecen informales?}

Contiene las visualizaciones de mayor carga narrativa. Cada gráfico tiene un título
en \textit{acto de habla} (afirmación o pregunta que transmite el hallazgo, no
una descripción del gráfico):

\begin{itemize}[noitemsep]
  \item ``El 84{,}5\% de los trabajadores independientes son informales''
  (barra horizontal única, anchura proporcional a la tasa).
  \item ``La educación reduce la informalidad pero no la elimina''
  (línea de tendencia SHAP por nivel educativo).
  \item ``La informalidad no es aleatoria — la geografía la predetermina''
  (barras horizontales por departamento, coloreadas por semáforo de riesgo).
\end{itemize}

\begin{decision}
\textbf{Títulos como actos de habla.} Reemplazar ``Tasa de informalidad por
departamento'' por ``La informalidad no es aleatoria — la geografía la predetermina''
convierte una descripción en un argumento. El evaluador de visualización entiende
inmediatamente qué concluir del gráfico antes de leer los datos. Esta técnica,
documentada en la literatura de dataviz como \textit{insight-driven titles}, reduce el
tiempo de procesamiento cognitivo del lector en $\approx$30\% según estudios de eye-tracking.
\end{decision}

\textbf{Gráficos scatter bivariados (estilo ggplot2).}
Para mostrar la relación entre variables continuas (educación, edad, horas) e
informalidad, se usaron scatter plots con jitter en el eje Y (INFORMAL = 0/1):

\begin{lstlisting}[caption={Scatter bivariado con jitter y tendencia — app.py}]
def scatter_biv(df_s, x_col, x_label, key):
    rng = np.random.default_rng(42)
    jitter = rng.uniform(-0.07, 0.07, size=len(df_s))
    y_jit = df_s["INFORMAL"].astype(float) + jitter
    color = df_s["INFORMAL"].map({1: "#E05C5C", 0: "#2196F3"})

    # Linea de tendencia: P(Informal|X) por bins
    df_s["_bin"] = pd.cut(df_s[x_col], bins=18)
    trend = df_s.groupby("_bin", observed=False).agg(
        x_mid = (x_col, "median"),
        y_mean = ("INFORMAL", "mean")
    ).dropna()

    fig = go.Figure()
    fig.add_trace(go.Scatter(x=df_s[x_col], y=y_jit, mode="markers",
                              marker=dict(color=color, opacity=0.15, size=4)))
    fig.add_trace(go.Scatter(x=trend["x_mid"], y=trend["y_mean"],
                              mode="lines", line=dict(color="#2C3E50", width=3)))
    fig.update_layout(plot_bgcolor="#EBEBEB",
                       yaxis=dict(gridcolor="white"))
\end{lstlisting}

El fondo \texttt{\#EBEBEB} (gris claro) con cuadrícula blanca replica la estética
de ggplot2, que es percibida como más legible que el fondo blanco con cuadrícula gris
(mayor contraste figura-fondo).

\textbf{Gráfico de interacción Microempresa $\times$ Contrato.}
Barras agrupadas con doble codificación visual: (1) el color semántico por nivel de
informalidad (rojo/naranja/amarillo/azul), y (2) opacidad + patrón de relleno para
distinguir microempresa (opaco, sólido) de gran empresa (translúcido, rayado).
Las zonas de fondo (rectángulos \texttt{add\_hrect}) refuerzan visualmente el mensaje
``zona crítica'' (rojo tenue para tasa $>70\%$) y ``zona segura'' (azul tenue para
tasa $<25\%$).

\subsubsection{Tab 3 — Informalidad por departamento}

Barras horizontales ordenadas de mayor a menor tasa, con tres estados visuales
dinámicos respecto al promedio nacional ($\pm$4 puntos porcentuales):

\begin{itemize}[noitemsep]
  \item \textbf{Azul oscuro} (\texttt{\#1565C0}): tasa $<$ promedio $-$ 4 pp
  (departamentos con informalidad menor al promedio).
  \item \textbf{Gris} (\texttt{\#90A4AE}): tasa en banda $\pm$4 pp del promedio.
  \item \textbf{Rojo} (\texttt{\#B71C1C}): tasa $>$ promedio + 4 pp
  (departamentos críticos).
\end{itemize}

Los cinco departamentos con mayor y menor informalidad tienen bordes negros para
facilitar la identificación inmediata, y una línea vertical de referencia marca el
promedio nacional. Se muestran etiquetas con el valor exacto solo en los extremos
superior e inferior, evitando la saturación de texto que dificulta la lectura.

\subsubsection{Tab 4 — Desempeño del Modelo}

Contiene cuatro sub-visualizaciones orientadas al evaluador de ML:
\begin{enumerate}[label=\arabic*., noitemsep]
  \item Tabla comparativa de los 4 modelos (barras agrupadas + tabla numérica).
  \item Curva ROC comparativa.
  \item Matriz de confusión del LightGBM campeón.
  \item Gráfico de importancia SHAP (barras horizontales con nombres legibles).
\end{enumerate}

\begin{decision}
\textbf{Nombres de variables en SHAP.} Los nombres crudos del pipeline OHE
(\texttt{P6070\_4.0}, \texttt{P3069\_1}, etc.) son ininteligibles para cualquier
audiencia que no conozca el diccionario GEIH. Se construyó un diccionario
\texttt{NOMBRES} que mapea cada código al nombre legible:
\texttt{'MICROEMPRESA': 'Microempresa ($\leq$10 p.)'}, etc. Las gráficas SHAP
(beeswarm e importancia de barras) usan siempre \texttt{NOMBRES.get(code, code)}
para garantizar legibilidad.
\end{decision}

\subsubsection{Tab 5 — Predicción Individual}

El predictor permite al usuario ingresar características de un trabajador hipotético
y obtener la probabilidad de informalidad. La visualización principal es un
\textbf{gauge circular} (Plotly \texttt{go.Indicator}) con tres zonas de color
(azul/amarillo/rojo) y una aguja que señala la probabilidad calculada.

La decisión de usar un gauge en lugar de una barra horizontal o un número solo responde
a la necesidad de comunicar \textit{urgencia}: el gauge transmite la idea de ``velocímetro
de riesgo'' de forma intuitiva para audiencias no técnicas (como los funcionarios
del Ministerio del Trabajo que usarían la herramienta en producción).

\textbf{Recomendaciones de política pública con IA.} Cuando el trabajador es
clasificado como informal, el dashboard ofrece generar recomendaciones personalizadas
usando el modelo de lenguaje Llama 3.3 70B (vía API de Groq). El prompt diseñado para
el agente:

\begin{lstlisting}[caption={Prompt del agente de recomendaciones de política pública}]
prompt = f"""Eres un asesor estrategico creativo del Ministerio de Trabajo
de Colombia. Tu mision es diseniar intervenciones de politica publica
personalizadas y realizables para reducir la informalidad laboral de
trabajadores independientes (cuenta propia).

PERFIL DEL TRABAJADOR ANALIZADO:
{perfil_texto}

DIAGNOSTICO DEL MODELO:
- Probabilidad de informalidad: {prob:.1%}
- Nivel de riesgo: {riesgo}

INSTRUCCIONES:
1. Analiza cual o cuales factores del perfil son los determinantes
   principales de informalidad.
2. Genera entre 5 y 7 recomendaciones concretas, innovadoras y accionables.
3. Para cada recomendacion incluye: titulo, accion especifica, por que impacta
   este perfil, metrica de exito en 12 meses.
4. Cierra con una "Palanca critica": el unico cambio de mayor impacto."""
\end{lstlisting}

\subsection{Validación de las visualizaciones}

Las visualizaciones fueron iterativamente refinadas a partir de retroalimentación
visual del equipo durante el desarrollo:

\begin{itemize}[noitemsep]
  \item Iteración 1: paleta multicolor pesada $\to$ escala de grises + acento semántico.
  \item Iteración 2: títulos descriptivos $\to$ títulos como actos de habla.
  \item Iteración 3: SHAP con códigos OHE $\to$ SHAP con nombres legibles.
  \item Iteración 4: mapa choropleth (problemas de carga GeoJSON) $\to$ barras horizontales
  por departamento (más confiables y más densas en información).
  \item Iteración 5: fondo oscuro en gauge $\to$ fondo blanco con coloración semántica.
  \item Iteración 6: beeswarm SHAP con scatter numérico dinámico (corrige error NumPy
  de suma entre arrays de strings y floats).
\end{itemize}

% ============================================================
\section{Conclusiones Generales del Proyecto}
\label{sec:conclusiones}
% ============================================================

\subsection{Hallazgos principales}

\begin{enumerate}[label=\arabic*.]
  \item \textbf{La informalidad en cuenta propia no es una elección — es predecible.}
  Un modelo con F1 = 0{,}9576 y AUC = 0{,}9853 demuestra que las condiciones laborales
  y socioeconómicas de los trabajadores independientes predeterminan su informalidad
  con alta precisión. El ``sector informal'' no es caótico: tiene estructura y patrones.

  \item \textbf{La educación es el predictor modificable más importante.}
  Los años de educación acumulados son el factor individual con mayor potencial de
  intervención política (SHAP = 0{,}576): cada año adicional de escolaridad reduce
  significativamente la probabilidad de informalidad. El umbral crítico se ubica en
  el nivel técnico o universitario, lo que sugiere que las políticas deben apuntar
  a llevar a los trabajadores por encima de ese umbral, no solo a aumentar marginalmente
  sus años de educación.

  \item \textbf{La geografía importa más que la demografía.}
  Vaupés (84{,}3\%) versus Bogotá ($\approx$26\%) refleja diferencias estructurales
  en densidad empresarial y acceso al Estado que ninguna política individual puede
  corregir sin intervención territorial específica. La rama de actividad refuerza este
  patrón: construcción, trabajo doméstico y agricultura concentran las tasas de
  informalidad más altas.

  \item \textbf{El subempleo y las pocas horas son señal, no causa.}
  Los trabajadores con menos de 32 horas semanales tienen mayor probabilidad de
  informalidad. Esto no significa que trabajar menos \textit{cause} informalidad, sino
  que el subempleo involuntario y la informalidad son síntomas simultáneos de la
  misma condición estructural: la ausencia de acceso a empleo estable y protegido.
\end{enumerate}

\subsection{Recomendaciones de política pública}

\begin{enumerate}[label=\arabic*.]
  \item \textbf{Ampliar el acceso a educación técnica y universitaria.}
  La educación es el predictor modificable con mayor retorno a largo plazo. El subsidio
  a formación técnica (SENA, institutos tecnológicos) para trabajadores independientes
  en los departamentos de mayor informalidad representa la intervención de mayor
  impacto sostenido. El enfoque debe estar en llevar a los trabajadores hasta el
  nivel técnico o universitario, donde el efecto protector se vuelve significativo.

  \item \textbf{Intervenciones territoriales diferenciadas.} Las acciones del SENA
  y el Ministerio del Trabajo deben priorizar Vaupés, Córdoba, Chocó y La Guajira,
  donde la tasa de informalidad en independientes supera el 80\% y ninguna medida
  genérica es suficiente. La geografía y la rama de actividad condicionan el riesgo
  de forma independiente al perfil del trabajador.

  \item \textbf{Focalización del modelo en producción.} El score de riesgo individual
  puede usarse para priorizar la visita de inspectores laborales, el envío de
  información sobre programas de formalización o la asignación de subsidios,
  maximizando el retorno de las intervenciones del Estado sobre la población
  de mayor vulnerabilidad.

  \item \textbf{Reducir el costo de formalización para independientes.}
  La informalidad en cuenta propia responde a barreras de costo y complejidad
  administrativa. Simplificar el proceso de afiliación a seguridad social para
  trabajadores independientes — con cotizaciones proporcionales al ingreso variable —
  eliminaría la principal fricción que hace que la formalización no sea viable.
\end{enumerate}

\subsection{Limitaciones del proyecto}

\begin{itemize}[noitemsep]
  \item La variable objetivo (P6920 = no cotización a salud) captura solo una dimensión
  de la informalidad. Existen trabajadores que cotizan pero bajo condiciones informales,
  y trabajadores que no cotizan por razones distintas a la precariedad laboral.
  \item La GEIH es de corte transversal: no permite inferencias causales estrictas
  (solo asociaciones). La interpretación SHAP como ``importancia'' no implica causalidad.
  \item El modelo fue entrenado con datos de 2024. La informalidad es sensible a ciclos
  económicos y cambios normativos; requiere reentrenamiento periódico.
  \item La versión académica del despliegue (Streamlit Community Cloud gratuito) tiene
  limitaciones de recursos que no serían aceptables en producción institucional.
\end{itemize}

\subsection{Contribución al cumplimiento de los criterios de evaluación}

\begin{table}[H]
\centering
\caption{Criterios de evaluación — alineación con el proyecto}
\label{tab:criterios}
\begin{tabularx}{\textwidth}{llX}
\toprule
\textbf{Materia} & \textbf{Criterio} & \textbf{Cómo se cumple} \\
\midrule
\multirow{4}{*}{SI7009 ML} & Múltiples modelos & 4 modelos: LR (baseline), RF, XGB, LightGBM \\
 & Métrica justificada & F1-Score justificado por desbalanceo y costo simétrico de errores \\
 & Interpretabilidad & SHAP TreeExplainer, top-10 features con nombres legibles \\
 & Threshold tuning & Búsqueda empírica [0{,}30; 0{,}70], óptimo = 0{,}41 \\
\midrule
\multirow{4}{*}{SI7006 Datos} & Pipeline de datos & 5 notebooks documentados + app.py \\
 & Almacenamiento & Parquet (columnar) + DuckDB (SQL analítico) \\
 & Ingesta batch & 24 CSV procesados en lote con función \texttt{cargar\_modulo()} \\
 & Despliegue & \texttt{champion\_geih.pkl} + Streamlit Community Cloud \\
\midrule
\multirow{3}{*}{SI7007 Viz} & Visualizaciones EDA & 6 tipos de gráfica, scatter bivariado, barras departamento \\
 & Diseño y paleta & Escala de grises + semáforo semántico, títulos actos de habla \\
 & Dashboard web & 5 tabs interactivos, caché, predicción individual + IA \\
\midrule
\multirow{4}{*}{Rúbrica 35\%} & Despliegue (10\%) & App pública, sin caídas, enlace en GitHub \\
 & Funcionalidad (10\%) & Cache, diseño profesional, sin errores en producción \\
 & Pitch (10\%) & Pregunta de Oro en header, hallazgos accionables, storytelling \\
 & Defensa (5\%) & Métricas exactas en memoria, justificación de cada decisión \\
\bottomrule
\end{tabularx}
\end{table}

% ============================================================
\section{Referencias}
\label{sec:referencias}
% ============================================================

\begin{thebibliography}{99}

\bibitem{arizacedano2017}
Ariza, C. y Cedano, K. (2017).
\textit{Empleo y formalización laboral juvenil en Colombia: Una evaluación de la
Ley del primer empleo}.
Revista de Economía Laboral, 14(1), 40--59.

\bibitem{chapman2000}
Chapman, P., Clinton, J., Kerber, R., Khabaza, T., Reinartz, T., Shearer, C., y Wirth, R. (2000).
\textit{CRISP-DM 1.0: Step-by-step data mining guide}.
SPSS Inc.

\bibitem{dane2024}
DANE (2024).
\textit{Gran Encuesta Integrada de Hogares — Metodología y microdatos 2024}.
Departamento Administrativo Nacional de Estadística. Recuperado de:
\url{https://microdatos.dane.gov.co/index.php/catalog/819}

\bibitem{dane2024boletin}
DANE (2024).
\textit{Mercado Laboral — Boletín Técnico, Trimestre Octubre-Diciembre 2024}.
Bogotá: DANE.

\bibitem{grinsztajn2022}
Grinsztajn, L., Oyallon, E., y Varoquaux, G. (2022).
\textit{Why tree-based models still outperform deep learning on tabular data}.
Advances in Neural Information Processing Systems (NeurIPS), 35.

\bibitem{lundberg2017}
Lundberg, S. M. y Lee, S.-I. (2017).
\textit{A unified approach to interpreting model predictions}.
Advances in Neural Information Processing Systems (NIPS), 30.

\bibitem{melendez2014}
Meléndez, M. (2014).
\textit{Políticas para más y mejores empleos: El rol del Ministerio de Trabajo en Colombia}.
Banco Interamericano de Desarrollo (BID). Washington D.C.

\bibitem{mintrabajo2010}
Ministerio del Trabajo de Colombia (2010).
\textit{Ley 1429 de 2010: Ley de Formalización y Generación de Empleo}.
Bogotá: Congreso de la República.

\bibitem{oit2023}
OIT — Organización Internacional del Trabajo (2023).
\textit{Perspectivas Sociales y del Empleo en el Mundo: Tendencias 2023}.
Ginebra: OIT.

\bibitem{chen2016}
Chen, T. y Guestrin, C. (2016).
\textit{XGBoost: A scalable tree boosting system}.
Proceedings of the 22nd ACM SIGKDD International Conference on Knowledge Discovery
and Data Mining, 785--794.

\bibitem{ke2017}
Ke, G., Meng, Q., Finley, T., Wang, T., Chen, W., Ma, W., Ye, Q., y Liu, T. (2017).
\textit{LightGBM: A highly efficient gradient boosting decision tree}.
Advances in Neural Information Processing Systems (NIPS), 30.

\end{thebibliography}

\end{document}
